% ========== 2.6 USE CASES ==========
% Descriptive scenarios showing system usage

\section{2.8 Use Cases}
\label{sec:use-cases}
\addcontentsline{toc}{section}{2.8 Use Cases}

\lettrine{U}{se cases} illustrate how different stakeholders interact with Symphony to accomplish their development goals. These scenarios demonstrate the system's capabilities across diverse contexts and user types, showing how AI-first architecture enables workflows that would be difficult or impossible with traditional development environments.

\subsection{Rapid Application Development}
\label{subsec:rapid-application-development}

A freelance developer receives a client request to build a web application for managing customer appointments. The developer opens Symphony and activates Maestro Mode, then describes the application requirements in natural language. The system analyzes the intent and activates the seven-model pipeline to transform the high-level description into working software.

The Enhancer-Prompt Model refines the vague requirements into detailed technical specifications, identifying necessary features like user authentication, appointment scheduling, email notifications, and administrative dashboards. The Feature Model breaks down these specifications into structured user stories with acceptance criteria. The Planner Model designs the application architecture, selecting appropriate technologies and frameworks based on the requirements.

The Coordinator Model sequences the implementation tasks, determining optimal order based on dependencies and risk. The Code-Visualizer Model creates pseudocode representations of key algorithms and workflows, helping the developer understand the implementation approach. The Editor Model generates production-ready source code for each component, following best practices and maintaining consistency. The Summarizer Model creates documentation explaining the system architecture, API endpoints, and deployment procedures.

Throughout this process, the developer maintains creative control, reviewing AI decisions at key checkpoints and providing guidance when the AI needs clarification. The entire application development, which might traditionally require several days of work, is completed in a few hours through effective human-AI collaboration.

\subsection{Legacy Code Modernization}
\label{subsec:legacy-modernization}

A development team inherits a legacy application written in outdated technologies that requires modernization. The codebase uses deprecated frameworks, lacks automated tests, and contains numerous code quality issues. The team lead opens the legacy project in Symphony and switches to Virtuoso Mode for context-sensitive refactoring assistance.

Symphony analyzes the codebase structure and identifies architectural patterns, dependencies, and potential modernization challenges. The system suggests a phased modernization approach that minimizes risk by updating components incrementally rather than attempting a complete rewrite. The team reviews this strategy and agrees with the proposed approach.

The team begins with the data access layer, where Symphony identifies database queries that could be optimized or replaced with modern patterns. The AI suggests specific refactoring operations that improve performance while maintaining behavioral compatibility. As the team works through different modules, Symphony learns their coding preferences and architectural decisions, adapting its suggestions to match the team's style.

The AI identifies security vulnerabilities in the legacy code, including injection risks and inadequate input validation. It suggests specific fixes that address these issues while explaining the security principles involved. Throughout the modernization process, Symphony maintains comprehensive documentation of changes, architectural decisions, and rationale. The project that initially seemed overwhelming becomes manageable through systematic AI-assisted refactoring.

\subsection{Collaborative Feature Development}
\label{subsec:collaborative-development}

A medium-sized development team works on adding a complex feature to an existing application. The feature requires coordination across frontend, backend, and database components, with multiple developers working simultaneously on different aspects. The team uses Symphony's collaborative capabilities to maintain consistency and avoid conflicts.

The technical lead creates a Melody workflow that defines the feature development process. This visual workflow specifies the sequence of tasks, dependencies between components, and quality gates that must be passed before integration. Team members can see the overall process and understand how their individual contributions fit into the larger feature.

Frontend developers work in Virtuoso Mode, building user interface components with AI assistance. Symphony suggests component structures that follow the team's established patterns and provides real-time feedback about accessibility and responsive design. Backend developers use Maestro Mode to implement API endpoints and business logic, describing the required functionality in natural language. Database developers receive AI assistance in designing schema changes and writing migration scripts.

As team members complete their components, Symphony's Artifact Store maintains version history and tracks dependencies between artifacts. When one developer modifies a shared interface, Symphony notifies other team members whose code depends on that interface. The team lead monitors progress through Symphony's analytics dashboard, which shows completion status, code quality metrics, and potential bottlenecks. The feature development that might traditionally require several weeks of coordination is completed efficiently through AI-assisted collaboration.

\subsection{Learning and Skill Development}
\label{subsec:learning-development}

A junior developer joins a team working with technologies and frameworks they have limited experience with. The developer uses Symphony's educational features to accelerate learning while contributing productively to the project. Symphony operates in Solo Mode initially, providing quick answers to specific questions without overwhelming the developer with complex orchestration.

When the developer encounters unfamiliar code patterns, they ask Symphony to explain the purpose and behavior. The AI provides clear explanations that connect the specific code to broader programming concepts, helping the developer understand not just what the code does, but why it was written that way. These explanations include references to relevant documentation and best practices.

As the developer gains confidence, they transition to Virtuoso Mode for more substantial coding tasks. Symphony provides suggestions that follow the team's established patterns, helping the developer write code that matches the team's style and standards. When the developer makes mistakes, Symphony explains the issues and suggests corrections, turning errors into learning opportunities.

The AI notices patterns in the developer's questions and proactively offers relevant information. When the developer repeatedly asks about asynchronous programming concepts, Symphony provides a mini-tutorial explaining promises, async-await syntax, and common pitfalls. This contextual education accelerates learning by addressing knowledge gaps as they become relevant to actual work.

Senior team members review the junior developer's code and provide feedback. Symphony incorporates this feedback into its assistance model, learning the team's specific preferences and standards. Future suggestions reflect this accumulated knowledge, helping the junior developer internalize team practices more quickly than would be possible through code review alone.

Over time, the developer's reliance on AI assistance decreases as their expertise grows. Symphony adapts to this progression, offering less basic assistance and more advanced suggestions that challenge the developer to consider edge cases, performance implications, and architectural concerns. The system serves as a patient mentor that adjusts its teaching approach based on the learner's growing capabilities.

\subsection{Enterprise Compliance and Governance}
\label{subsec:enterprise-compliance}

A large financial services organization adopts Symphony for application development while maintaining strict compliance with regulatory requirements. The organization's security team configures Symphony with enterprise-grade controls that enforce security policies without impeding developer productivity.

The security team deploys Symphony with custom extensions that implement organization-specific security checks. These extensions scan all code for potential vulnerabilities, enforce secure coding standards, and prevent accidental exposure of sensitive data. Developers receive immediate feedback when their code violates security policies, with clear explanations of the issues and suggested remediation.

All AI interactions are logged to a centralized audit system that maintains comprehensive records for compliance verification. When regulators request evidence of security controls, the organization can demonstrate that all code generation and modification activities were monitored and complied with established policies. The audit logs include details about what code was generated, which AI models were used, and what security checks were applied.

The organization maintains a private extension repository containing proprietary AI models trained on internal codebases and business logic. These models understand organization-specific patterns and requirements that public AI models cannot address. Developers benefit from AI assistance that reflects institutional knowledge while maintaining confidentiality of proprietary information.

Access controls ensure that developers can only use AI features appropriate for their roles and clearance levels. Junior developers receive more restrictive AI assistance that prevents potentially dangerous operations, while senior developers have broader autonomy. This role-based access control balances productivity with risk management.

The organization's architecture team creates approved Melody workflows that encode best practices and compliance requirements. Developers use these workflows for common tasks, ensuring consistent adherence to organizational standards. The workflows include automated quality gates that prevent non-compliant code from reaching production systems.

Performance monitoring reveals that Symphony's AI assistance reduces the time developers spend on routine coding tasks, allowing them to focus on complex business logic and architectural decisions. The organization measures productivity improvements and code quality enhancements, demonstrating clear return on investment that justifies the adoption costs.

\subsection{Open Source Contribution}
\label{subsec:open-source-contribution}

A developer wants to contribute to an open-source project but finds the large codebase intimidating. They use Symphony to understand the project structure and identify appropriate areas for contribution. Symphony analyzes the codebase and provides a high-level overview of the architecture, explaining how different components interact and where various features are implemented.

The developer identifies a bug report in the project's issue tracker and wants to fix it. They describe the issue to Symphony, which analyzes the relevant code and suggests potential causes. The AI explains the code's intended behavior and identifies where the actual behavior diverges from expectations. This analysis helps the developer understand the bug without spending hours manually tracing through unfamiliar code.

Symphony suggests a fix that addresses the bug while maintaining compatibility with the rest of the codebase. The AI explains why this particular approach is appropriate, considering factors like the project's coding style, architectural patterns, and backward compatibility requirements. The developer reviews the suggestion, makes minor adjustments based on their understanding of the issue, and implements the fix.

Before submitting the contribution, the developer uses Symphony to generate comprehensive tests that verify the fix works correctly and doesn't introduce regressions. The AI creates test cases that cover edge cases the developer might not have considered, improving the quality of the contribution. Symphony also generates documentation explaining the fix and updating relevant comments in the code.

The developer submits a pull request that includes the fix, tests, and documentation. The project maintainers review the contribution and request some modifications to better align with project conventions. The developer uses Symphony to understand the requested changes and implement them quickly. The contribution is accepted, and the developer has successfully participated in open-source development with AI assistance that made the process less intimidating and more productive.

This use case demonstrates how Symphony democratizes participation in complex software projects by reducing the barriers to understanding and contributing to unfamiliar codebases. The AI serves as a knowledgeable guide that helps developers navigate complexity and produce high-quality contributions.
